
% ============================================================
% FRENCH TENSES — PASSÉ COMPOSÉ
% Version graphique LaTeX/TikZ
% Compilation : XeLaTeX
% ============================================================
\documentclass[a4paper,10pt]{article}

\usepackage[a4paper,top=9mm,bottom=8mm,left=11mm,right=11mm]{geometry}
\usepackage{fontspec}
\usepackage[french]{babel}
\usepackage{microtype}
\usepackage{tikz}
\usetikzlibrary{arrows.meta,calc,positioning}
\usepackage[most]{tcolorbox}
\usepackage{tabularx}
\usepackage{array}
\usepackage{xcolor}
\usepackage[normalem]{ulem}

\setmainfont{Noto Sans}
\setsansfont{Noto Sans}

% ------------------------------------------------------------
% PALETTE — PASSÉ = BLEU
% ------------------------------------------------------------
\definecolor{Blue}{HTML}{356FA8}
\definecolor{BlueDark}{HTML}{244D78}
\definecolor{BlueLight}{HTML}{EAF2F9}
\definecolor{BlueLine}{HTML}{9BBBD7}
\definecolor{Border}{HTML}{B9D3E3}

\definecolor{Red}{HTML}{D52F35}
\definecolor{RedLight}{HTML}{FCEBED}
\definecolor{Gold}{HTML}{B27A00}
\definecolor{GoldLight}{HTML}{FFF3D7}
\definecolor{Ink}{HTML}{172A32}
\definecolor{Grey}{HTML}{64757D}
\definecolor{White}{HTML}{FFFFFF}

\pagestyle{empty}
\setlength{\parindent}{0pt}
\setlength{\parskip}{0pt}
\setlength{\tabcolsep}{3pt}

% ------------------------------------------------------------
% CADRE
% ------------------------------------------------------------
\newcommand{\PageFrame}{%
\begin{tikzpicture}[remember picture,overlay]
  \coordinate (A) at ([xshift=8mm,yshift=-8mm]current page.north west);
  \coordinate (B) at ([xshift=-8mm,yshift=-8mm]current page.north east);
  \coordinate (C) at ([xshift=-8mm,yshift=8mm]current page.south east);
  \coordinate (D) at ([xshift=8mm,yshift=8mm]current page.south west);
  \def\r{1.7mm}

  \draw[Border,line width=.65pt]
    (A)
    -- ([xshift=-\r]B)
    arc[start angle=90,end angle=0,radius=\r]
    -- ([yshift=\r]C)
    arc[start angle=0,end angle=-90,radius=\r]
    -- (D) -- cycle;

  \draw[Blue,line width=3.2mm] (A) -- (D);
\end{tikzpicture}%
}

% ------------------------------------------------------------
% BOÎTES
% ------------------------------------------------------------
\tcbset{
  enhanced,
  boxrule=.6pt,
  arc=2.8mm,
  left=3.7mm,right=3.7mm,top=1.75mm,bottom=1.75mm,
  boxsep=1mm,
  coltext=Ink
}

\newtcolorbox{bluebox}{
  colback=BlueLight,colframe=BlueLine
}
\newtcolorbox{redbox}{
  colback=RedLight,colframe=Border
}
\newtcolorbox{whitebox}{
  colback=White,colframe=Border
}
\newtcolorbox{goldbox}{
  colback=GoldLight,colframe=GoldLight,
  boxrule=0pt,arc=2.5mm,
  left=2.7mm,right=2.7mm,top=1.25mm,bottom=1.25mm
}

\newcommand{\sectiontitle}[2]{%
  {\large\bfseries\textcolor{#1}{#2}}%
}

\newcommand{\pill}[2][BlueLight]{%
  \tikz[baseline=(X.base)]{
    \node[
      fill=#1,rounded corners=2pt,
      inner xsep=4.2pt,inner ysep=1.6pt,
      font=\small\bfseries,text=Ink
    ] (X) {#2};
  }%
}

\newcommand{\signalpill}[1]{%
  \tikz[baseline=(X.base)]{
    \node[
      fill=GoldLight,rounded corners=2pt,
      inner xsep=3.7pt,inner ysep=1.1pt,
      font=\scriptsize\bfseries,text=Ink
    ] (X) {#1};
  }%
}

% ------------------------------------------------------------
% TIMELINE — ACTION TERMINÉE DANS LE PASSÉ, LIÉE AU PRÉSENT
% ------------------------------------------------------------
\newcommand{\Timeline}{%
\begin{tikzpicture}[x=1cm,y=1cm]
  \draw[-{Latex[length=3mm]},very thick,Ink] (0,0) -- (15.3,0);

  % NOW
  \draw[-{Latex[length=3mm,width=2mm]},very thick,Blue]
    (9.3,1.10) -- (9.3,.10);
  \node[font=\bfseries\small,text=Blue] at (9.3,1.32) {MAINTENANT};

  % Past completed action
  \fill[Blue] (5.1,0) circle (3.2pt);
  \draw[Blue,line width=2.2pt] (5.1,0) -- (9.0,0);

  % subtle link to present
  \draw[Blue!45,dashed,line width=.7pt] (5.1,0.18) -- (9.0,0.18);

  \node[
    fill=BlueLight,rounded corners=2pt,
    inner xsep=5pt,inner ysep=2pt,
    font=\scriptsize\bfseries,text=BlueDark
  ] at (6.9,.72) {action terminée / lien avec le présent};

  \node[font=\scriptsize\bfseries,text=Grey] at (2.2,-.48) {PASSÉ};
  \node[font=\scriptsize\bfseries,text=Blue] at (9.3,-.48) {PRÉSENT};
  \node[font=\scriptsize\bfseries,text=Grey] at (13.3,-.48) {FUTUR};
\end{tikzpicture}%
}

\begin{document}
\PageFrame

% ------------------------------------------------------------
% TITRE
% ------------------------------------------------------------
\vspace*{1mm}
\begin{center}
  {\fontsize{25}{27}\selectfont\bfseries\textcolor{Ink}{PASSÉ COMPOSÉ}}

  \vspace{0.75mm}
  {\large\textcolor{Grey}{Le passé composé}}

  \vspace{3mm}

  \begin{tcolorbox}[
    width=.88\textwidth,
    colback=BlueLight,
    colframe=BlueLight,
    boxrule=0pt,
    arc=3mm,
    halign=center,
    top=2.2mm,bottom=2.2mm
  ]
    {\large\bfseries\textcolor{BlueDark}{AUXILIAIRE AU PRÉSENT + PARTICIPE PASSÉ}}
    \quad
    {\small(\textbf{avoir} ou \textbf{être})}
  \end{tcolorbox}
\end{center}

\vspace{1mm}

% ------------------------------------------------------------
% TIMELINE
% ------------------------------------------------------------
\begin{bluebox}
  \sectiontitle{Blue}{La ligne du temps}

  \vspace{0.75mm}
  \begin{center}
    \Timeline
  \end{center}

  \vspace{-1mm}
  \begin{center}
    \small
    Le \textbf{passé composé} exprime surtout une action
    \textbf{\textcolor{Blue}{terminée}} dans le passé, souvent avec
    \textbf{\textcolor{Blue}{un lien avec le présent}}.
  \end{center}
\end{bluebox}

\vspace{1.35mm}

% ------------------------------------------------------------
% CONSTRUCTION
% ------------------------------------------------------------
\begin{bluebox}
  \sectiontitle{Blue}{Construction}

  \vspace{0.75mm}
  \begin{tabularx}{\textwidth}{@{}p{.29\textwidth} X@{}}
    \textbf{Avec avoir} &
      J'\textbf{ai parlé}. \qquad Nous \textbf{avons fini}.\\
    \textbf{Avec être} &
      Elle \textbf{est partie}. \qquad Ils \textbf{sont arrivés}.
  \end{tabularx}

  \vspace{0.7mm}
  \textit{La majorité des verbes se conjuguent avec \textbf{avoir}.}
\end{bluebox}

\vspace{1.35mm}

% ------------------------------------------------------------
% FORME / ACCORD
% ------------------------------------------------------------
\begin{minipage}[t]{.485\textwidth}
\begin{bluebox}
  \sectiontitle{Blue}{Exemples}

  \vspace{0.75mm}
  J'\textbf{ai travaillé}.\\[-1mm]
  Tu \textbf{as compris}.\\[-1mm]
  Nous \textbf{avons mangé}.\\[-1mm]
  Elle \textbf{est arrivée}.
\end{bluebox}
\end{minipage}
\hfill
\begin{minipage}[t]{.485\textwidth}
\begin{redbox}
  \sectiontitle{Red}{Attention !}

  \vspace{0.75mm}
  Avec \textbf{être}, le participe passé
  s'accorde avec le sujet :

  \vspace{0.4mm}
  Elle est arriv\textbf{ée}.\\[-1mm]
  Ils sont arriv\textbf{és}.
\end{redbox}
\end{minipage}

\vspace{1.35mm}

% ------------------------------------------------------------
% USES
% ------------------------------------------------------------
\begin{bluebox}
  \sectiontitle{Blue}{Quand utilise-t-on le passé composé ?}

  \vspace{0.75mm}
  \begin{tabularx}{\textwidth}{@{}p{.31\textwidth} X@{}}
    \pill{ACTION TERMINÉE}
      &
      Action achevée dans le passé.
      \newline \textit{J'ai terminé mon travail.}\\[0.7mm]

    \pill{ÉVÉNEMENT PONCTUEL}
      &
      Événement qui s'est produit à un moment donné.
      \newline \textit{Il est arrivé à huit heures.}\\[0.7mm]

    \pill{BILAN / EXPÉRIENCE}
      &
      Expérience ou bilan en rapport avec le présent.
      \newline \textit{J'ai déjà vu ce film.}\\[0.7mm]

    \pill{SUCCESSION D'ACTIONS}
      &
      Actions successives dans un récit.
      \newline \textit{Il est entré, a posé son sac et s'est assis.}
  \end{tabularx}
\end{bluebox}

\vspace{1.35mm}

% ------------------------------------------------------------
% COMPARISON
% ------------------------------------------------------------
\begin{whitebox}
  \sectiontitle{Blue}{Passé composé ou imparfait ?}

  \vspace{0.65mm}
  \begin{tabularx}{\textwidth}{@{}p{.45\textwidth} X@{}}
    \textbf{\textcolor{Blue}{Passé composé}} &
      \textbf{\textcolor{Grey}{Imparfait}}\\
    action terminée, événement &
      description, habitude, durée en cours\\[-.2mm]
    \textit{J'ai mangé.} &
      \textit{Je mangeais.}
  \end{tabularx}
\end{whitebox}

\vspace{1.35mm}

% ------------------------------------------------------------
% SIGNAL WORDS
% ------------------------------------------------------------
\begin{goldbox}
  \begin{center}
    {\normalsize\bfseries\textcolor{Gold}{REPÈRES FRÉQUENTS}}
    \hspace{3mm}
    \signalpill{hier}\hspace{1mm}
    \signalpill{déjà}\hspace{1mm}
    \signalpill{jamais}\hspace{1mm}
    \signalpill{encore}

    \vspace{0.75mm}

    \signalpill{récemment}\hspace{1mm}
    \signalpill{ce matin}\hspace{1mm}
    \signalpill{la semaine dernière}\hspace{1mm}
    \signalpill{tout à l'heure}
  \end{center}
\end{goldbox}

\end{document}
